\documentclass{article}
\usepackage{amsmath,amsthm}
\newtheorem{theorem}{Theorem}
\begin{document}
\title{Transporting a reserve check across changes of membership}
\maketitle

\paragraph{Scope.}
Fix a service cut and an additive obligation already assigned to a site.
Retained service supply is counted only if available at that cut under
effective permission. A newcomer can replace the old supply lease, inherit an
old exposure, or change a provenance declaration; roster entry alone is not a
grant or a physical effect. This model checks a fixed obligation. It does not
synthesize a new controller or bind an unknown failure law.

\paragraph{Composition and a reusable checked cut.}
Let $r$ be a retained state, $e$ an incoming continuation, and $b$ the
shared provenance, timing and obligation values needed to compose them.
Write $a(r)$ for retained capacity and $x(e),d(e)$ for extra capacity and
amount owed. Let $m_b$ be the minimum of $a(r)$ over retained states at $b$,
with an empty cell treated separately. In the finite model such a minimum
has a witness. At fixed ports, a module may change the extension histories
without rebuilding the old minimum profile.

\begin{theorem}
If the join allows all matched pairs for each retained slice used by the
incoming module, the universal service check is exactly
$m_{b(e)}+x(e)\ge d(e)$ for all admitted $e$. The same conclusion about
exactness holds if each admitted $e$ is compatible with at least one
old minimizing state in its slice. Regardless of these compatibility
assumptions, the scalar inequality is sufficient when all actually
admitted retained states were already included in the lower bound.
\end{theorem}

The argument uses the realized minimizer for necessity and monotonicity
of addition for sufficiency. A slice that the old model did not admit has
no old minimum. If new wiring excludes all old minimizers behind the same
boundary, a failed inequality is not evidence of a real shortfall. For
example, feasible pairs of retained and incoming capacity might be only
$(0,1)$ and $(1,0)$, while the relaxed join also admits $(0,0)$.
Include the relevant coupling in the interface or check the explicitly
allowed relation. For elementary consumers that can test a requested
capacity on a single slice, two different minima are distinguishable
by requesting the larger minimum. This gives minimal sufficient
information for those consumers, not a memory bound for every physical
controller.

\paragraph{Civil development comparisons.}
One fixture counts staging units before pickup; another counts report
forwarding opportunities at retry. Incidental old work changes retained
supply without changing the named domain. The nominal outage law permits
at most one named domain failure. The old guarantee still holds with a
timely new member on a distinct domain: when the old domain is down the
new lease supplies the obligation, and when the new domain is down retained
capacity suffices. With the otherwise identical new lease on the old
domain, one outage suffices to miss the obligation. The new member's
individual possible availability and the age and amount of its loan
need not differ. A non-overlapping domain does not rescue a lease that
becomes active after the cut.

The frozen certificate is computed once for each retained fixture.
The exact joint comparator generates every outage and incidental-work
combination from the same declared law. A strong incremental comparator
caching the same slice minima gives an operational tie. A coarse summary
using only the minimum of all the old states cannot validate both feasible
and infeasible replacements at existing information. This is a verified
extension pattern with authored plant adapters, not independently selected
mission acceptance or a measured fault law. Old plant support, provenance
and identity-to-effect bindings are adapter obligations; the checker cannot
infer them from a roster.
\end{document}
